% !PRINTED none
% Section-opening page (no running-head folio visible in any of the three bands):
% large upper blank area, then centred "Fifth Section." / "The Modular Functions."
% title block with a rule, followed directly by the "§ 51." heading and the
% start of its body text on the same leaf.
\begin{center}
\gesperrt{Fifth Section.}

\textbf{The Modular Functions.}

\rule{2cm}{0.4pt}

\textbf{\S\,51. \gesperrt{The Elliptic Differential Equations.}}
\end{center}

% The source's closing predicate, omitted in the recovered German fragment, has been restored.
The elliptic functions defined so far are functions of the two variables $v$, $\omega$, and the moduli $\varkappa$, $\varkappa'$, and further $j(\omega)$, $K$, $K'$ depend on the ratio of periods $\omega$, which must have a positive imaginary part. Likewise the function $\wp(u)$ is a function of the three variables $u$, $\omega_1$, $\omega_2$, and the invariants $g_2$, $g_3$ depend on $\omega_1$, $\omega_2$. \gesperrt{The problem of the inversion of the elliptic integrals}, as we formulated it in \S\,14, presupposes however that in the elliptic functions $\varkappa^2$ (or, in the \gesperrt{Weierstrass} normal form, $g_2, g_3$) have arbitrarily given values, and the complete solution of this problem therefore requires that not $\omega$, but $\varkappa^2$ (respectively $g_2, g_3$) appear as the second independent variable, and that $\omega$ be expressed by these. The next considerations will be devoted to this problem.

We start from the problem of integrating a system of differential equations, which we call that of the \gesperrt{elliptic differential equations}:
\weq{1}{\begin{aligned}
\frac{dx}{dv} &= yz, \\
\frac{dy}{dv} &= -zx, \\
\frac{dz}{dv} &= -\varkappa^2xy,
\end{aligned}}
under the subsidiary conditions:
\weq{2}{\text{for } v=0 \text{ we are to have } x=0,\ y=1,\ z=1\text{.}}
